\documentclass{beamer}
\usetheme{default}
\usecolortheme{seahorse}

\usepackage{amsmath,amssymb,mathrsfs,bm}
\usepackage{anyfontsize}
\usepackage{unicode-math}
\unimathsetup{math-style=ISO, bold-style=ISO, mathrm=sym}
\setmathfont{XITS Math}
\usepackage{fontspec}
\setsansfont{Libertinus Sans}

\AtBeginDocument{
  \let\uglyepsilon\epsilon
  \renewcommand\epsilon{\varepsilon}
}
\AtBeginDocument{
  \let\uglymathsf\mathsf
  \renewcommand\mathsf{\symsf}
}
\renewcommand{\bm}{\symbf}

\usepackage{color}
\usepackage{graphicx,caption,subcaption}
\usepackage{hyperref}
\hypersetup{
    colorlinks=true,
    linkcolor=blue,
    filecolor=green,
    urlcolor=cyan,
    citecolor=cyan,
}

\title{Kiloparsec-scale AGN Outflows and Feedback in Merger Free Galaxies}
\author{Zhijun Zhang}
\date{2021.10.28}

\begin{document}

\begin{frame}[plain]
    \maketitle
\end{frame}

\begin{frame}{Introduction}
    \textbf{Recent observations and simulations have challenged the long-held paradigm that mergers are the dominant mechanism driving the growth of both galaxies and supermassive black holes.}
    \begin{itemize}
        \item[Old] Galaxy growth mainly occurs through mergers
        \item[New] Galaxy growth mainly occurs through non-merger “secular processes” rather than by mergers
    \end{itemize}
    \begin{itemize}
        \item[\bullet] Millenium simulations: the majority of bulge mass is built through disk instabilities, Parry et al. (2009)
        \item[\bullet] Observations: only 27\% of star formation is triggered by mergers at $z\sim 2$, Kaviraj et al. (2013)
    \end{itemize}
\end{frame}

\begin{frame}{Observations}
    A significant bulge component still forms even in a minor merger

    Select 4 disk-dominated galaxies with luminous, unobscured Type 1 AGN at $z<0.1$
    \begin{figure}
    \includegraphics[width=0.6\linewidth]{Figure_1.jpg}
    \includegraphics[width=0.6\linewidth]{Figure_2.jpg}
    \end{figure}
    All four galaxies have little or no bulge component.

    Harry and Theodore show a strong bar feature, Neville a weak bar feature, and Padma a barlens feature.
\end{frame}

\begin{frame}{Observations}
    \begin{itemize}
        \item[\bullet] \textcolor{blue}{narrow H$\mathrm{\beta}$} shows the galaxy structure
        \item[\bullet] \textcolor{blue}{narrow $\mathrm{O_{III}}$} ionized by the central AGN
        \item[\bullet] \textcolor{blue}{broad $\mathrm{O_{III}}$} ionized by the AGN outflow
    \end{itemize}
\end{frame}

\begin{frame}{Observations}
    \begin{columns}
        \begin{column}{0.5\linewidth}
            \begin{figure}
                \includegraphics[width=1\linewidth]{Figure_3.png}
            \end{figure}
        \end{column}
        \begin{column}{0.5\linewidth}
            \textbf{Flux Density - Wavelength}

            \begin{itemize}
                \item[\textcolor{black}{black}] \textcolor{black}{Spectrum}
                \item[\textcolor{red}{red}] \textcolor{red}{Fittings}
                \item[\textcolor{gray}{cross}] \textcolor{gray}{Residuals}
            \end{itemize}

            \begin{itemize}
                \item[$\mathrm{O_{III}}$]
                \begin{itemize}
                    \item[\bullet] \textcolor{blue}{narrow}
                    \item[\bullet] \textcolor{green}{broad}
                \end{itemize}
                \item[H$\mathrm{\beta}$]
                \begin{itemize}
                    \item[\bullet] \textcolor{blue}{narrow}
                    \item[\bullet] \textcolor{green}{broad}
                    \item[\bullet] \textcolor{magenta}{broad line region}
                \end{itemize}
            \end{itemize}
        \end{column}
    \end{columns}
\end{frame}

\begin{frame}{Observations}
    \begin{figure}
        \includegraphics[width=1\linewidth]{Figure_4.jpg}
        \caption*{\textcolor{blue}{narrow H$\mathrm{\beta}$} integrated flux, velocity, velocity dispersion}
    \end{figure}
\end{frame}

\begin{frame}{Observations}
    \begin{figure}
        \includegraphics[width=1\linewidth]{Figure_5.jpg}
        \caption*{\textcolor{blue}{narrow $\mathrm{O_{III}}$} integrated flux, velocity, velocity dispersion}
    \end{figure}
\end{frame}

\begin{frame}{Observations}
    \begin{figure}
        \includegraphics[width=1\linewidth]{Figure_6.jpg}
        \caption*{\textcolor{blue}{broad $\mathrm{O_{III}}$} integrated flux, velocity, velocity dispersion}
    \end{figure}
\end{frame}

\begin{frame}{Observations}
    \begin{figure}
        \includegraphics[width=1\linewidth]{Figure_7_1.jpg}
        \caption*{\textcolor{blue}{narrow $\mathrm{O_{III}/H\beta}$ emission}}
    \end{figure}
    \begin{figure}
        \includegraphics[width=1\linewidth]{Figure_7_2.jpg}
        \caption*{\textcolor{blue}{broad $\mathrm{O_{III}/H\beta}$ emission}}
    \end{figure}
\end{frame}

\begin{frame}{Calculations}
    \begin{itemize}
        \item[\bullet] Outflow luminosity from broad $\mathrm{O_{III}}$ gives gas mass in the outflow:
        \begin{small}
        \[
            M_{\mathrm{O_{III}}} = 8\times10^7M_{\astrosun}\left(\frac{\overline{n_e}^2/\overline{n_e^2}}{10^{\mathrm{[O/H]-[O/H]}_\astrosun}}\right)\left(\frac{L[\mathrm{O_{III}}]}{10^{44}\,\mathrm{erg/s}}\right)\left(\frac{n_e}{500\,\mathrm{cm^{-3}}}\right)^{-1}
        \]
        \end{small}
        assuming that $n_e=500\,\mathrm{cm^{-3}}$ and $\mathrm{[O/H]=[O/H]}_\astrosun$
        \item[\bullet] Maximum velocity of the outflow:
        \[
            v_{\mathrm{O_{III}}} = |\mathrm{\Delta}v_{max}| + 2\sigma_{broad,max}
        \]
    \end{itemize}
\end{frame}

\begin{frame}{Calculations}
    \begin{itemize}
        \item[\bullet] Physical extent of the outflow $r_{max}$ (deconvolved by seeing)
        \item[\bullet] Timescale of the outflow $t_{\mathrm{out}} = \dfrac{r_{max}}{v_{\mathrm{O_{III}}}}$
        \item[\bullet] Mass, kinetic energy and momentum outflow rate:
        \[
            \dot{M} = B\frac{M_{\mathrm{O_{III}}}}{t_{\mathrm{out}}},\quad
            \dot{E} = \frac{1}{2}\dot{M}v_{\mathrm{O_{III}}}^2,\quad
            \dot{P} = \dot{M}v_{\mathrm{O_{III}}}
        \]
        assuming geometry factor $B=1$
        \item[\bullet] Black hole accretion rates $\dot{m}=\eta L_{\mathrm{bol}}c^2$ where the radiative efficiency $\eta=0.15$ and $L_{\mathrm{bol}}$ using the W3 $12\,\mathrm{\mu m}$ band
    \end{itemize}
\end{frame}

\begin{frame}{Results}
    \begin{figure}
        \includegraphics[width=1\linewidth]{Table_2.png}
    \end{figure}
    \begin{itemize}
        \item[\bullet] $\overline{M}_{\mathrm{O_{III}}}=5.5\pm0.2M_\astrosun$
        \item[\bullet] $\overline{\dot{M}}=0.3\pm0.1\,\mathrm{M_\astrosun/yr}$
        \item[\bullet] $\overline{r_{max}}=1.6\pm0.4\,\mathrm{kpc}$, $\sim25\%$ of the galaxy Petrosian radius
    \end{itemize}
\end{frame}

\begin{frame}{Discussion}
    \begin{itemize}
        \item[\bullet] Bars and spiral arms would be capable of driving inflows which could sustain both the SMBH growth ($0.02\sim0.07\,\mathrm{M_\astrosun/yr}$) and an outflow of AGN ($0.2\sim0.7\,\mathrm{M_\astrosun/yr}$) for an extended period of $0.6\sim1.9\,\mathrm{Myr}$

        \item[\bullet] Secular processes are responsible for the majority of SBMH growth, whereas mergers are responsible for the majority of outflowing material and the subsequent feedback on the galaxy

        \item[\bullet] The average ratio $\log_{10}(c\dot{P}/L_{\mathrm{bol}})=-0.91\pm0.08$ suggests that these outflows are momentum-conserving

        \item[\bullet] The average $v_{\mathrm{O_{III}}}$ of the 4 targets is $\sim30.5$ times larger than escape velocity, which will cause AGN feedback
    \end{itemize}
\end{frame}

\begin{frame}{Discussion}
    \begin{itemize}
        \item[\bullet] The outflow rates in Bae et al. (2017) are $\sim51$ times larger than in this article, but SMBH accretion rates are $\sim3$ times lower, this may cause by a spin up of SMBHs due to a secular feeding mechanism

        \item[\bullet] A higher spatial resolution IFU study and a larger sample of disk-dominated galaxies would allow for a more detailed study
    \end{itemize}
\end{frame}

\begin{frame}{Conclusion}
    \begin{itemize}
        \item[\bullet] Secularly powered outflows are typical of low-redshift Type 1 AGNs and they have momentum conserving outflows

        \item[\bullet] Secular processes are responsible for the majority of SMBH growth over cosmic time
    \end{itemize}
\end{frame}

\begin{frame}
    \begin{center}
        \Huge Thank you for listening!
    \end{center}
\end{frame}

\end{document}
